\section*{Chapitre \ref{ch:iv1:stochastic-calculus} --- Calcul stochastique}
\markboth{\protect\omTocMark{}{Réponses~: Calcul stochastique}}{\protect\omTocMark{}{Réponses~: Calcul stochastique}}

\begin{solution}{iq:iv1:stochastic-calculus:1}
$\operatorname{Cov}(W_s, W_t) = \min(s, t)$. La variance vaut 5, la corrélation $\frac12$ et la probabilité $\frac38$.

Pour $s \le t$, on écrit $W_t = W_s + (W_t - W_s)$, dont le second terme est indépendant de $W_s$. Il reste $\operatorname{Var}W_s = s$.

Alors $\operatorname{Var}(W_1 + W_2) = 1 + 2 + 2(1) = 5$. La corrélation de $W_1$ et $W_4$ vaut $1/\sqrt{1 \times 4} = \frac12$.

Pour deux variables gaussiennes conjointes, centrées et de corrélation $r$, la probabilité que les deux soient positives vaut $\frac14 + \arcsin(r)/(2\pi)$\index{probabilité d'orthant}. Ici $r = 1/\sqrt{2}$, donc $\arcsin r = \pi/4$, et la probabilité vaut $\frac14 + \frac18$. Une simulation donne $0.375$.

\iqwrong{Traiter $W_1$ et $W_2$ comme indépendants, ce qui donne une variance de 3 et une probabilité de $\frac14$. Ils se recouvrent sur $[0, 1]$.}

\iqfollow{Quelle est la variance de $W_1 + \dots + W_{10}$~?}{$\sum_{i, j}\min(i, j) = 1^2 + 2^2 + \dots + 10^2 = 385$, une identité utile.}

\iqwhy{La formule de covariance est la source de tout le reste, et la question vérifie qu'elle peut être utilisée et pas seulement récitée.}
\end{solution}

\begin{solution}{iq:iv1:stochastic-calculus:2}
Martingales~: $W_t$, $W_t^2 - t$, $e^{W_t - t/2}$ et $W_t^3 - 3tW_t$. Ne sont pas des martingales~: $W_t^2$, $tW_t$ et $W_t^3$. Le test consiste à savoir si $f_t + \frac12f_{xx} = 0$.

D'après le lemme d'Itô\index{lemme d'Itô}, $df(t, W_t) = (f_t + \frac12f_{xx})\,dt + f_x\,dW$. Une martingale exige que le terme en $dt$ s'annule.

$W_t$ donne $0$. $W_t^2$ donne $1$, donc il a une dérive, et retrancher $t$ l'élimine. $tW_t$ donne $W_t$, qui n'est pas nul. $e^{W_t - t/2}$ donne $-\frac12 f + \frac12f = 0$. $W_t^3$ donne $3W_t$, et $W_t^3 - 3tW_t$ donne $-3W_t + 3W_t = 0$. Une quadrature gaussienne des espérances conditionnelles confirme les sept verdicts.

Une règle empirique pour $tW_t$~: sa moyenne vaut 0, ce qui peut suggérer une martingale. Une moyenne constante est nécessaire, mais pas suffisante. Son espérance conditionnelle sachant le temps $s$ vaut $tW_s$, et non $sW_s$.

\iqwrong{Affirmer que $W_t^3$ est une martingale parce que sa moyenne est toujours nulle. Sachant $W_s$, la moyenne de $W_t^3$ vaut $W_s^3 + 3(t - s)W_s$.}

\iqfollow{Trouvez le polynôme suivant de la famille.}{$W_t^4 - 6tW_t^2 + 3t^2$. Ce sont des polynômes d'Hermite en $W_t/\sqrt{t}$, multipliés par des puissances de $\sqrt{t}$.}

\iqwhy{Le test en une ligne remplace sept calculs séparés. Le connaître témoigne d'une maîtrise pratique du lemme d'Itô.}
\end{solution}

\begin{solution}{iq:iv1:stochastic-calculus:3}
$d\ln S = (\mu - \frac12\sigma^2)\,dt + \sigma\,dW$. La moyenne vaut $e^{0.1} = 1.105$, la médiane $e^{0.08} = 1.083$, et la probabilité $0.46$.

Le lemme d'Itô donne $d\ln S = dS/S - \frac12(dS)^2/S^2$, et $(dS)^2 = \sigma^2S^2\,dt$. Le logarithme a donc une dérive $\mu - \frac12\sigma^2 = 8\%$\index{mouvement brownien géométrique}.

$\ln S_1$ suit une loi normale de moyenne $0.08$ et d'écart type $0.2$. La médiane est à la moyenne du logarithme, donc $e^{0.08}$. Pour la moyenne de $S$, $\mathbb{E}e^{0.08 + 0.2Z} = e^{0.08 + 0.02} = e^{0.1}$, comme le montre aussi l'équation pour $S$.

La probabilité de finir au-dessus de la moyenne revient à $0.08 + 0.2Z > 0.1$, soit $Z > 0.1$, avec probabilité $1 - \Phi(0.1) = 0.46$. Une simulation de l'équation discrétisée donne ces valeurs.

Sans volatilité, la moyenne et la médiane coïncident. L'écart entre les deux croît avec $\sigma^2$.

\iqwrong{Omettre le $-\frac12\sigma^2$, qui est le terme d'Itô d'ordre deux, et donner une médiane de $e^{0.1}$.}

\iqfollow{À quelle volatilité la croissance médiane est-elle nulle, avec $\mu = 10~\%$~?}{$\sigma^2/2 = 0.1$, donc $\sigma = 44.7~\%$. La trajectoire moyenne croît toujours de 10~\% alors que la trajectoire typique ne va nulle part.}

\iqwhy{C'est la première application du lemme d'Itô en finance, et l'écart entre moyenne et médiane en est le résultat le plus utile.}
\end{solution}

\begin{solution}{iq:iv1:stochastic-calculus:4}
$\mathbb{E}W_t^4 = 3t^2$, $\mathbb{E}e^{aW_t} = e^{a^2t/2}$, $\mathbb{E}e^{2W_1} = e^2 \approx 7.39$, et la corrélation vaut $\frac14$.

Puisque $W_t = \sqrt{t}Z$, le moment d'ordre quatre vaut $t^2\mathbb{E}Z^4 = 3t^2$. Compléter le carré dans l'intégrale gaussienne donne la formule exponentielle\index{fonction génératrice des moments}.

Pour les carrés, on écrit $W_4 = W_1 + V$, où $V$ est indépendant de $W_1$ et de variance 3. Alors
\[
\mathbb{E}W_1^2W_4^2 = \mathbb{E}W_1^4 + \mathbb{E}W_1^2\,\mathbb{E}V^2 = 3 + 3 = 6.
\]
Donc la covariance vaut $6 - 1 \times 4 = 2$. Les variances valent $2$ et $2 \times 16 = 32$, ce qui donne $2/\sqrt{64} = \frac14$. C'est le carré de la corrélation $\frac12$ de $W_1$ et $W_4$, ce qui est toujours vrai pour deux variables gaussiennes conjointes centrées. Une simulation confirme ce résultat.

\iqwrong{Donner $\mathbb{E}W_t^4 = t^2$, ce qui néglige le kurtosis. Ou donner la corrélation $\frac12$, celle de $W_1$ et $W_4$ elles-mêmes.}

\iqfollow{Quelle est $\mathbb{E}W_t^6$~?}{$15t^3$. Les moments pairs de la loi normale valent $1, 3, 15, 105$, chacun étant le précédent multiplié par le nombre impair suivant.}

\iqwhy{Ces moments reviennent dans chaque calcul de variance. Les connaître par cœur fait gagner du temps, et la corrélation des carrés est un petit test de rigueur.}
\end{solution}

\enlargethispage{-3\baselineskip}%
\begin{solution}{iq:iv1:stochastic-calculus:5}
$P(W_4 > 2) = P(W_1 > 1) = 0.1587$, et la volatilité annuelle vaut $15.9~\%$.

$X_t = cW_{t/c^2}$ est continu, part de 0 et a des accroissements indépendants de loi normale centrée. Sa covariance vaut $c^2\min(s/c^2, t/c^2) = \min(s, t)$\index{changement d'échelle brownien}. C'est donc un mouvement brownien.

Avec $c = 2$~: $W_4 = 2W'_1$ pour un mouvement brownien $W'$, donc $P(W_4 > 2) = P(W'_1 > 1)$, soit $1 - \Phi(1) = 0.1587$.

Des rendements indépendants ajoutent leurs variances, donc les écarts types croissent comme la racine carrée du temps~: $1~\% \times \sqrt{252} = 15.9~\%$. Si les rendements ne sont pas indépendants, ce chiffre est faux, et c'est pourquoi une mesure d'autocorrélation doit l'accompagner.

\iqwrong{Multiplier par 252, ce qui donne 252~\%. La variance croît avec le temps, et la volatilité avec sa racine.}

\iqfollow{Quelle est la volatilité annuelle pour une volatilité mensuelle de 5~\%~?}{$5~\% \times \sqrt{12} = 17.3~\%$.}

\iqwhy{Le changement d'échelle est la propriété brownienne la plus utilisée sur n'importe quel desk, et la démontrer avec la formule de covariance prend une ligne.}
\end{solution}

\begin{solution}{iq:iv1:stochastic-calculus:6}
La somme des carrés a une moyenne de 1 et un écart type de $0.045$. La somme des valeurs absolues a une moyenne de $25.2$. Quand les pas rétrécissent, la première converge vers 1, et la seconde diverge.

Chaque accroissement suit une loi normale de variance $1/n$. Son carré a pour moyenne $1/n$ et pour variance $2/n^2$. En sommant sur $n$ pas indépendants, on obtient une moyenne de 1 et une variance de $2/n$\index{variation quadratique}, donc la somme converge vers 1 et son écart type pour $n = 1000$ vaut $\sqrt{0.002}$.

La valeur absolue a pour moyenne $\sqrt{2/\pi}/\sqrt{n}$, donc la somme a pour moyenne $\sqrt{2n/\pi}$, soit $25.2$ ici. Elle croît comme $\sqrt{n}$\index{variation totale}. Une trajectoire brownienne est donc de longueur infinie mais de variation quadratique finie, et c'est la raison d'être du terme supplémentaire du lemme d'Itô~: $(dW)^2$ se comporte comme $dt$.

Une simulation de $20\,000$ trajectoires confirme les trois valeurs.

\iqwrong{Penser que la somme des carrés reste aléatoire à la limite, ou qu'elle tend vers 0 comme pour une trajectoire régulière. Pour une trajectoire régulière, la somme des carrés tend vers 0, alors qu'ici elle reste égale à $T$, sans aléa résiduel.}

\iqfollow{À quel point une variance réalisée est-elle bruitée avec 252 rendements quotidiens~?}{L'écart type relatif vaut $\sqrt{2/252} = 8.9~\%$ pour une volatilité constante et des rendements gaussiens, et davantage avec des queues épaisses.}

\iqwhy{C'est le cœur du sujet. Cela explique $(dW)^2 = dt$ par un calcul, et c'est la théorie sous-jacente à la variance réalisée.}
\end{solution}

\begin{solution}{iq:iv1:stochastic-calculus:7}
$\int_0^TW\,dW = \frac12(W_T^2 - T)$, de moyenne 0 et de variance $T^2/2$. La somme à gauche converge vers cette valeur, la somme à droite vers cette valeur plus $T$, et le point milieu vers $\frac12W_T^2$.

Le lemme d'Itô appliqué à $W^2$ donne $d(W^2) = 2W\,dW + dt$, dont on tire le résultat\index{intégrale d'Itô}. La moyenne vaut 0, puisque $\mathbb{E}W_T^2 = T$. La variance vaut $\frac14\operatorname{Var}W_T^2 = \frac14 \times 2T^2$, ou, par l'isométrie d'Itô\index{isométrie d'Itô}, $\int_0^T\mathbb{E}W_t^2\,dt = T^2/2$.

La somme à gauche est la définition, et ses termes sont de moyenne nulle, car l'intégrande est connu avant l'accroissement.

La somme de droite moins la somme de gauche vaut $\sum(\Delta W)^2$, qui converge vers $T$, de sorte que la somme de droite a pour moyenne $T$. À la limite, elle vaut $\frac12(W_T^2 + T)$.

La moyenne des deux sommes est la forme au point milieu, qui se télescope exactement~: $\sum\frac12(W_i + W_{i + 1})(W_{i + 1} - W_i) = \frac12W_T^2$. C'est l'intégrale de Stratonovich\index{intégrale de Stratonovich}, qui conserve les règles du calcul ordinaire, et évaluer $W$ au véritable milieu de chaque pas conduit à la même limite. Une simulation à 1000 pas montre les sommes à gauche, à droite et moyennée.

\iqwrong{Écrire $\frac12W_T^2$ à partir du calcul ordinaire. Sa moyenne vaut $T/2$, ce qui est impossible pour une somme de termes de moyenne nulle.}

\iqfollow{Pourquoi la finance utilise-t-elle l'extrémité gauche~?}{Parce qu'une position doit être choisie avant le mouvement de prix qu'elle rapporte. Une règle à l'extrémité droite reviendrait à trader sur le futur, et elle gagnerait exactement le $T$ supplémentaire.}

\iqwhy{C'est l'illustration la plus courte de la raison pour laquelle les intégrales stochastiques diffèrent des intégrales ordinaires, et elle montre que le choix de l'extrémité a un contenu réel.}
\end{solution}

\enlargethispage{-1\baselineskip}%
\begin{solution}{iq:iv1:stochastic-calculus:8}
$I_T$ est normale de moyenne 0 et de variance $T^3/3$, soit 9 pour $T = 3$. Sa corrélation avec $W_T$ vaut $\sqrt{3}/2 = 0.866$.

En intégrant par parties\index{mouvement brownien intégré}, $d(tW_t) = W_t\,dt + t\,dW_t$, d'où
\[
I_T = TW_T - \int_0^Tt\,dW_t = \int_0^T(T - t)\,dW_t.
\]
C'est une intégrale d'une fonction déterministe contre $dW$, elle est donc normale de moyenne 0. La variance, par l'isométrie, vaut $\int_0^T(T - t)^2\,dt = T^3/3$.

Pour vérifier, l'intégrale double donne $\int_0^T\int_0^T\min(s, t)\,ds\,dt = T^3/3$, et pour $T = 3$ une quadrature donne également 9.

La covariance avec $W_T = \int_0^T dW$ vaut $\int_0^T(T - t)\,dt = T^2/2$. La corrélation vaut donc
\[
\frac{T^2/2}{\sqrt{T^3/3 \cdot T}} = \frac{\sqrt{3}}{2},
\]
quelle que soit la valeur de $T$. Une simulation à 2000 pas confirme ce résultat.

\iqwrong{Donner la variance sous la forme $T^2/2$ ou $T^3$. Une vérification des unités aide~: $I_T$ a pour unité celle de $W$ multipliée par un temps, donc sa variance croît comme $T \times T^2$.}

\iqfollow{Comment la volatilité du prix moyenné sur une période se compare-t-elle à celle du prix lui-même~?}{La moyenne $I_T/T$ a pour variance $T/3$, donc sa volatilité est celle du prix divisée par $\sqrt{3}$. Une volatilité de 20~\% devient $11.5~\%$, ce qui explique pourquoi les options asiatiques sont moins chères.}

\iqwhy{Le mouvement brownien intégré est l'un des calculs standard, et l'intégration par parties est plus éclairante que l'intégrale double.}
\end{solution}

\begin{solution}{iq:iv1:stochastic-calculus:9}
$P(M_T \ge a) = 2P(W_T \ge a)$, donc $M_T$ suit la loi de $|W_T|$. $P(M_1 \ge 1) = 0.317$, et $\mathbb{E}M_1 = 0.798$. Ne vérifier qu'aux clôtures quotidiennes déplace la barrière d'environ $0.5826\sigma\sqrt{\Delta}$.

Dès que $W$ atteint pour la première fois $a$, la suite de la trajectoire est symétrique par rapport à $a$\index{principe de réflexion}, si bien que la fin a autant de chances de se situer au-dessus de $a$ qu'en dessous. Donc $P(M_T \ge a, W_T < a) = P(M_T \ge a, W_T > a) = P(W_T > a)$. En additionnant, on obtient le facteur 2.

Ainsi $P(M_T \ge a) = 2(1 - \Phi(a/\sqrt{T}))$, qui est la queue de $|W_T|$. Pour $a = T = 1$, elle vaut $0.3173$. La moyenne vaut $\mathbb{E}|W_1| = \sqrt{2/\pi}$.

Une trajectoire peut traverser la barrière entre deux clôtures et revenir. Le maximum discret est donc plus bas, et la probabilité d'atteinte est plus faible. La correction usuelle\index{surveillance discrète} décale la barrière de $0.5826$ écart type d'un pas, soit ici $0.5826/\sqrt{252} = 0.037$. On obtient $2(1 - \Phi(1.037)) = 0.300$. Une simulation des clôtures quotidiennes sur $100\,000$ trajectoires donne ce nombre, et non $0.317$.

\iqwrong{Oublier le facteur 2, ou utiliser la formule continue pour une barrière contrôlée une fois par jour. Sur une barrière d'un an, cela représente environ deux points de probabilité.}

\iqfollow{Pourquoi le facteur 2 ne tient-il pas en présence de dérive~?}{La trajectoire réfléchie n'est plus aussi probable. Avec une dérive, la formule gagne un second terme, affecté d'un facteur $e^{2\mu a/\sigma^2}$.}

\iqwhy{C'est le résultat standard pour les options à barrière, et la correction pour la surveillance discrète est là où la théorie rencontre une véritable fiche de termes.}
\end{solution}

\begin{solution}{iq:iv1:stochastic-calculus:10}
À partir de la loi $P(T_a \le t) = 2(1 - \Phi(a/\sqrt{t}))$, la médiane pour $a = 1$ vaut $2.20$, et $P(T_1 > 100) = 0.080$. La moyenne est infinie, car la queue est d'ordre $1/\sqrt{t}$.

$T_a \le t$ est l'événement $M_t \ge a$\index{temps de premier passage}, donc le principe de réflexion donne la formule. Lorsque $t \to \infty$, elle tend vers $2(1 - \Phi(0)) = 1$, si bien que le temps est fini avec probabilité un.

Pour la médiane, $2(1 - \Phi(1/\sqrt{t})) = \frac12$ donne $1/\sqrt{t} = \Phi^{-1}(0.75) = 0.6745$, donc $t = 2.198$.

Par ailleurs, $P(T_1 > 100) = 1 - 2(1 - \Phi(0.1)) = 0.0797$.

Pour $t$ grand, $P(T_a > t) \approx 2\phi(0)a/\sqrt{t} = 0.798a/\sqrt{t}$. La moyenne est l'intégrale de cette queue, qui diverge. De manière équivalente, $T_a$ suit la loi de $a^2/Z^2$. Des échantillons de cette forme confirment à la fois la médiane et la queue, et les moyennes tronquées continuent de croître à mesure que le seuil augmente. Une exécution directe de la trajectoire donne un atteinte avant $t = 2.198$ environ une fois sur deux.

\iqwrong{Supposer qu'un événement certain doit avoir un temps d'attente moyen fini. Ici, la moitié des trajectoires arrivent avant $2.2$, tandis que 8~\% sont encore en route à 100.}

\iqfollow{Comment la médiane change-t-elle pour $a = 10$~?}{Par changement d'échelle, elle est multipliée par $a^2$, si bien que la médiane vaut environ 220.}

\iqwhy{Une arrivée certaine avec une moyenne infinie est la surprise principale du sujet, et elle réapparaît chaque fois qu'on suppose qu'un pari de moyenne nulle reviendra dans un temps connu.}
\end{solution}

\begin{solution}{iq:iv1:stochastic-calculus:11}
La probabilité de finir en 2 vaut $\frac13$, et le temps espéré vaut 2. Plus généralement, ces quantités valent $a/(a + b)$ et $ab$.

Soit $\tau$ le temps d'arrêt. Il est d'espérance finie, et le $W$ arrêté est borné, de sorte que le théorème d'arrêt optionnel\index{arrêt optionnel} s'applique.

Avec $W$~: $\mathbb{E}W_\tau = 0$. Si $q$ est la probabilité de finir en $b$, alors $qb - (1 - q)a = 0$, donc $q = a/(a + b)$. Avec $a = 1$ et $b = 2$, cela donne $\frac13$.

Avec $W^2 - t$~: $\mathbb{E}W_\tau^2 = \mathbb{E}\tau$. Donc
\[
\mathbb{E}\tau = qb^2 + (1 - q)a^2 = \frac{ab^2 + ba^2}{a + b} = ab = 2.
\]
Une marche symétrique sur réseau donne exactement ces mêmes valeurs, et une simulation sur $20\,000$ trajectoires sur un réseau de pas $0.1$ donne $0.33$ et $2.0$.

Lorsque $a \to \infty$, $q \to 1$ et $\mathbb{E}\tau = ab \to \infty$~: atteindre $b$ est certain mais prend un temps espéré infini, ce qui concorde avec le temps d'atteinte unilatéral.

\iqwrong{Utiliser le théorème d'arrêt optionnel sans vérification. Pour le problème unilatéral, il donnerait $\mathbb{E}W_\tau = 0$, alors que $W_\tau = b$. C'est la borne sur le processus arrêté qui le rend valide ici.}

\iqfollow{Quelles sont les réponses pour un intervalle symétrique $\pm a$~?}{$\frac12$ et $a^2$. Ainsi, le temps typique pour parcourir une distance $a$ est $a^2$, la même échelle que la variance.}

\iqwhy{Deux martingales et une ligne chacune. C'est l'usage canonique du théorème d'arrêt optionnel, et les recruteurs guettent la vérification de son domaine d'application.}
\end{solution}

\begin{solution}{iq:iv1:stochastic-calculus:12}
$e^{-2\mu X_t/\sigma^2}$ est une martingale, ce qui donne $P(\text{atteindre un jour } {-a}) = e^{-2\mu a/\sigma^2}$. Ici, $e^{-1} = 0.368$.

Essayons $e^{\theta X_t}$. Sa dérive a pour facteur $\theta\mu + \frac12\theta^2\sigma^2$, qui s'annule pour $\theta = -2\mu/\sigma^2$\index{martingale exponentielle}.

Arrêtons-nous en $\tau$, le premier instant où l'on atteint $-a$, ou à un niveau élevé $b$. La martingale est bornée sur l'intervalle, donc le théorème d'arrêt optionnel donne
\[
1 = qe^{2\mu a/\sigma^2} + (1 - q)e^{-2\mu b/\sigma^2},
\]
où $q$ est la probabilité d'atteindre $-a$ en premier. Lorsque $b \to \infty$, le second terme s'annule, donc $q = e^{-2\mu a/\sigma^2}$.

Avec $\mu = 0.1$, $\sigma = 0.2$ et $a = 0.2$~: $2 \times 0.1 \times 0.2/0.04 = 1$, donc la réponse est $e^{-1}$. Une marche ayant la même moyenne et la même variance vérifie exactement la formule de ruine, et elle donne $0.368$ pour de petits pas. Une simulation par pas de $0.01$ an confirme ce résultat, après la correction liée à la vérification à des instants discrets.

Remarquons les limites~: $\mu \to 0$ donne 1, car la trajectoire sans dérive atteint tout niveau, et un grand $\mu$ donne un nombre petit. Le rapport clé est $\mu a/\sigma^2$, donc doubler la volatilité divise l'exposant par quatre.

\iqwrong{Utiliser une queue gaussienne sur un seul horizon, par exemple une chance d'un an de $-20~\%$. Cela ne tient pas compte de toutes les trajectoires qui plongent puis remontent, et le chiffre obtenu est bien plus petit que la valeur réelle.}

\iqfollow{Quelle est la probabilité de perdre un jour 50~\%, avec les mêmes paramètres~?}{$e^{-2.5} = 0.082$. La probabilité décroît exponentiellement avec l'ampleur de la perte.}

\iqwhy{Il s'agit d'une formule pratique de risque de ruine. On la dérive en trois lignes à l'aide d'une martingale, et elle est souvent nécessaire dans une discussion sur le dimensionnement.}
\end{solution}

\begin{solution}{iq:iv1:stochastic-calculus:13}
$\mathbb{E}S_T^2 = S_0^2e^{(2\mu + \sigma^2)T}$, la puissance de martingale est $p = 1 - 2\mu/\sigma^2 = -4$, et la probabilité vaut $15/15.94 = 0.941$.

Le lemme d'Itô donne
\[
d(S^p) = S^p\left(\left(p\mu + \tfrac12p(p - 1)\sigma^2\right)dt + p\sigma\,dW\right).
\]
Ainsi $\mathbb{E}S_T^p = S_0^p\exp((p\mu + \frac12p(p - 1)\sigma^2)T)$, et $p = 2$ donne le moment d'ordre deux, $e^{0.24} = 1.271$ par unité de carré initial en $T = 1$.

La dérive est nulle lorsque $\mu + \frac12(p - 1)\sigma^2 = 0$, ou en $p = 0$, ce qui est trivial. Donc $p = 1 - 2\mu/\sigma^2 = 1 - 5 = -4$\index{fonction d'échelle}.

En arrêtant $S^{-4}$ en $2$ ou en $\frac12$, où elle est bornée~: $q2^{-4} + (1 - q)2^4 = 1$. Donc $q = 15/(16 - \frac{1}{16}) = 0.941$.

Une seconde voie travaille avec $\ln S$, de dérive $0.08$ et de volatilité $0.2$, avec sortie en $\pm\ln 2$. La martingale exponentielle $e^{\theta X}$ avec $\theta = -4$ donne la même équation, ce qui n'a rien de surprenant, puisque $S^{-4} = e^{-4\ln S}$. Une simulation à pas fins donne $0.94$.

\iqwrong{Affirmer que la probabilité vaut $\frac12$, puisque les deux niveaux sont à un facteur 2 l'un de l'autre, ou qu'elle découle de $\mu$ seul. La dérive logarithmique de $0.08$ face à une volatilité de $0.2$ est ce qui rend le résultat si déséquilibré.}

\iqfollow{Quelle est $p$ lorsque $\mu = \frac12\sigma^2$~?}{$p = 0$, et la martingale est alors le logarithme lui-même, puisque le logarithme n'a pas de dérive. La probabilité de sortie vaut alors exactement $\frac12$.}

\iqwhy{Trouver la puissance qui annule la dérive est la voie la plus rapide pour résoudre les problèmes de sortie pour les processus géométriques.}
\end{solution}

\begin{solution}{iq:iv1:stochastic-calculus:14}
L'espérance vaut $sx$, et la variance $s(1 - s)$. Sachant $W_1 = 0$, la probabilité d'atteindre 1 vaut $e^{-2} = 0.135$. En général, pour $a \ge \max(x, 0)$, elle vaut $e^{-2a(a - x)}$.

$(W_s, W_1)$ est conjointement gaussien, de covariance $s$ et $\operatorname{Var}W_1 = 1$. La régression donne l'espérance conditionnelle $sx$ et la variance $s - s^2$\index{pont brownien}. Le pont est donc fixé aux deux extrémités, avec l'incertitude maximale, $\frac14$, au milieu.

Par réflexion en $a$, une trajectoire qui finit en $x < a$ et a touché $a$ correspond à une trajectoire qui finit en $2a - x$. La densité des trajectoires ayant les deux propriétés est donc $\phi(2a - x)$, et la probabilité conditionnelle est le rapport des densités~:
\[
\frac{\phi(2a - x)}{\phi(x)} = \exp\left(-\frac{(2a - x)^2 - x^2}{2}\right) = e^{-2a(a - x)}.
\]
Avec $x = 0$ et $a = 1$, cela vaut $e^{-2}$. Remarquons la limite $x \to a$, où la formule donne 1, comme il se doit. Un pont simulé à 2000 pas donne $0.135$ après la correction liée à la vérification à des instants discrets.

\iqwrong{Utiliser $2(1 - \Phi(1))$ à partir du maximum non conditionnel. Cela moyenne sur toutes les extrémités. Fixer l'extrémité en 0 réduit la probabilité de plus de moitié.}

\iqfollow{Un cours passe de 100 à 101 en une journée, avec une volatilité journalière de 1. Quelle est la probabilité qu'il ait touché 102 entre les deux clôtures~?}{Ici $a = 2$ et $x = 1$, d'où $e^{-4} = 0.018$. C'est ainsi que le suivi de barrière est corrigé entre les dates d'observation.}

\iqwhy{Cette formule figure dans tout pricer de barrière et dans toute simulation dépendant de la trajectoire. La démonstration par réflexion est assez courte pour un entretien.}
\end{solution}

\begin{solution}{iq:iv1:stochastic-calculus:15}
$X_t = X_0e^{-\kappa t} + \sigma\int_0^te^{-\kappa(t - s)}\,dW_s$. L'espérance vaut $0.406$, la variance $0.245$, la variance stationnaire $\frac14$, et la demi-vie $0.347$.

Multiplier par le facteur intégrant\index{processus d'Ornstein-Uhlenbeck}~: $d(e^{\kappa t}X) = e^{\kappa t}\sigma\,dW$. En intégrant, on obtient la solution.

L'espérance vaut $X_0e^{-\kappa t}$, ici $3e^{-2} = 0.406$. Par l'isométrie, la variance vaut
\[
\sigma^2\int_0^te^{-2\kappa(t - s)}\,ds = \frac{\sigma^2}{2\kappa}\left(1 - e^{-2\kappa t}\right),
\]
qui vaut $\frac14(1 - e^{-4}) = 0.245$ pour ces valeurs. Lorsque $t$ croît, elle tend vers la valeur stationnaire $\sigma^2/(2\kappa) = \frac14$. La demi-vie est la solution de $e^{-\kappa t} = \frac12$, donc $t = \ln 2/\kappa$.

Limites~: pour $t$ petit, la variance est proche de $\sigma^2t$, comme pour le mouvement brownien standard. Lorsque $\kappa \to 0$, il en va de même pour tout $t$. Un échantillonnage exact AR(1) et une simulation d'Euler donnent $0.406$ et $0.245$.

\iqwrong{Oublier que le retour à la moyenne réduit aussi la contribution du bruit, et donner la variance comme $\sigma^2t$. Elle est bornée, et c'est justement l'intérêt du modèle.}

\iqfollow{Une série journalière a un coefficient AR(1) de $0.98$. Quels sont $\kappa$ et la demi-vie~?}{$\kappa = -\ln 0.98 = 0.0202$ par jour, et la demi-vie vaut $34.3$ jours. C'est le calcul standard pour un spread à retour à la moyenne.}

\iqwhy{C'est le deuxième processus stochastique le plus utilisé en finance, à la base de nombreux modèles de pairs trading et de taux.}
\end{solution}

\begin{solution}{iq:iv1:stochastic-calculus:16}
$W^2 = \rho W^1 + \sqrt{1 - \rho^2}\,W^\perp$, $\mathbb{E}W^1_tW^2_t = \rho t$, et les deux valeurs sont 3 et $1.5$.

La construction a une variance unitaire, et la corrélation avec $W^1$ est $\rho$\index{mouvements browniens corrélés}.

Le lemme d'Itô donne $d(W^1W^2) = W^1\,dW^2 + W^2\,dW^1 + \rho\,dt$, puisque $dW^1dW^2 = \rho\,dt$. Les intégrales stochastiques sont d'espérance nulle, donc l'espérance du produit vaut $\rho t$.

Ensuite, $\operatorname{Var}(W^1 + W^2) = t + t + 2\rho t = 3$.

Pour le dernier, $\mathbb{E}[X^2Y^2] = \sigma_X^2\sigma_Y^2 + 2\operatorname{Cov}(X, Y)^2$ pour des variables conjointement gaussiennes d'espérance nulle, d'après le théorème d'Isserlis\index{théorème d'Isserlis}. Cela donne $t^2(1 + 2\rho^2) = 1.5$. Une substitution directe le vérifie~: en écrivant $W^2 = \frac12W^1 + \sqrt{\frac34}V$ avec $V$ indépendant, on obtient $\frac14 \times 3 + \frac34 \times 1 = 1.5$. Une simulation le confirme. En $\rho = 0$, le résultat vaut 1, produit des variances, et en $\rho = 1$ il vaut 3, soit le moment d'ordre quatre.

\iqwrong{Multiplier les variances pour obtenir 1, ce qui ignore la corrélation des carrés. Ou écrire $W^2 = \rho W^1 + (1 - \rho)W^\perp$, qui n'est pas de variance unitaire.}

\iqfollow{Quelle est la variance de la somme pour $\rho = -\frac12$~?}{$2t(1 + \rho) = 1$ en $t = 1$. C'est la formule de base de la diversification.}

\iqwhy{Toute simulation multi-actifs se construit de cette façon. Le terme croisé de la formule d'Itô est l'endroit où l'on se trompe.}
\end{solution}

\begin{solution}{iq:iv1:stochastic-calculus:17}
L'identité découle du décalage de la densité. L'estimateur translaté réduit l'erreur type relative de 18~\% à environ $0.2~\%$, soit un facteur supérieur à 60.

Avec $Z = W_T$ de densité $\phi_T$,
\[
\mathbb{E}[f(W_T + \mu T)] = \int f(x + \mu T)\phi_T(x)\,dx = \int f(y)\phi_T(y - \mu T)\,dy,
\]
et en complétant le carré, on obtient $\phi_T(y - \mu T) = \phi_T(y)e^{\mu y - \mu^2T/2}$\index{théorème de Girsanov}. Il est donc possible de reporter la dérive dans un poids. C'est le théorème de Girsanov pour un seul instant, et il s'étend aux trajectoires.

À titre de vérification, avec $\mu = T = 1$ et $f = \mathbf{1}(x > 0)$, les deux formes donnent $\Phi(1) = 0.8413$.

Pour $P(W_1 > 4) = 3.17 \times 10^{-5}$, une simulation directe à $10^6$ trajectoires obtient environ 32 tirages dans l'événement, et a donc une erreur relative de $18~\%$.

Inversons l'identité~: simulons $Y = Z + 4$ et moyennons $\mathbf{1}(Y > 4)e^{-4Y + 8}$\index{échantillonnage préférentiel}. Désormais la moitié des trajectoires tombent dans l'événement, et chacune porte un poids faible. Sur une exécution, cela réduit l'erreur relative à environ $0.2\%$, ce qui concorde avec une intégrale pour le moment d'ordre deux. L'échantillonnage direct demanderait plusieurs milliers de fois plus de trajectoires pour la même précision.

Un décalage trop grand nuit à nouveau, car les poids deviennent erratiques. Pour un événement de queue, le bon choix place la moyenne à la frontière de l'événement.

\iqwrong{Oublier le poids, ou utiliser $e^{4Y - 8}$, dans le mauvais sens. Un contrôle vérifiant que le poids moyen vaut 1 révèle cette erreur aussitôt.}

\iqfollow{Pourquoi ce poids est-il une martingale en temps~?}{C'est la martingale exponentielle $e^{\mu W_t - \mu^2t/2}$, d'espérance 1, et c'est elle qui fait de la nouvelle mesure une mesure de probabilité.}

\iqwhy{Le changement de probabilité est souvent énoncé de façon abstraite. Cette version montre ce qu'il fait en une ligne, et donne une raison pratique de s'y intéresser.}
\end{solution}

\begin{solution}{iq:iv1:stochastic-calculus:18}
Probabilité réelle $0.655$, probabilité risque-neutre $0.460$, prix $0.460$. Le prix de marché du risque vaut $\frac12$. Le prix ne dépend pas de $\mu$.

$\ln S_1 - \ln 100$ est gaussien d'écart type $0.2$. Son espérance vaut $\mu - \frac12\sigma^2 = 0.08$ sous la probabilité réelle, donc $P(S_1 > 100) = \Phi(0.08/0.2) = \Phi(0.4)$.

Sous la probabilité risque-neutre, la dérive est le taux, ici nul\index{probabilité risque-neutre}, donc l'espérance du logarithme vaut $-0.02$, et la probabilité vaut $\Phi(-0.1) = 0.460$. Avec un taux nul, le prix est cette probabilité, par l'argument de réplication d'un digital.

Le prix de marché du risque vaut $\lambda = (\mu - r)/\sigma = 0.5$\index{prix de marché du risque}. Les deux probabilités sont liées par $dQ/dP = e^{-\lambda W_1 - \lambda^2/2}$, ce qui retire la dérive $\lambda\sigma$ du rendement. Repondérer les trajectoires réelles par ce facteur donne $0.460$ dans une simulation, et le poids est d'espérance 1.

Quiconque détient le digital à $0.46$ a une espérance de gain positive, d'un montant de $0.655 - 0.46$. Ce n'est pas gratuit~: le digital paie dans les états où l'action monte, il porte donc un risque sur l'action, dont le détenteur attend une rémunération.

\iqwrong{Fixer le prix avec la probabilité réelle, ou croire que $\mu$ influe sur le prix parce qu'il influe sur la chance du paiement. La couverture le retire du prix, mais pas de l'espérance de gain du détenteur.}

\iqfollow{Quelles sont les probabilités si $\mu = 2~\%$~?}{La chance réelle vaut désormais $\Phi(0) = \frac12$, et le prix reste $0.460$. L'écart s'est réduit à $0.04$, car l'espérance de gain ne vient que de la dérive.}

\iqwhy{La confusion entre probabilités et prix est l'erreur la plus fréquente dans ce domaine, et le digital l'isole.}
\end{solution}

\begin{solution}{iq:iv1:stochastic-calculus:19}
où $u(t, x) = e^{-r(T - t)}\mathbb{E}[g(X_T) \mid X_t = x]$, où $dX = \mu\,dt + \sigma\,dW$. Les deux solutions sont $x^2 + T - t$ et $e^{x + (T - t)/2}$.

Soit $M_s = e^{-rs}u(s, X_s)$. Le lemme d'Itô donne une dérive $e^{-rs}(u_t + \mu u_x + \frac12\sigma^2u_{xx} - ru)$, qui est nulle par l'équation\index{formule de Feynman-Kac}. Donc $M$ est une martingale, et $M_t = \mathbb{E}[M_T \mid X_t]$, ce qui est la formule.

Avec $\mu = r = 0$, $\sigma = 1$~: sachant $W_t = x$, $W_T$ est gaussien, d'espérance $x$ et de variance $T - t$.

Pour $g = x^2$~: $u = x^2 + (T - t)$. Vérification~: $u_t = -1$ et $\frac12u_{xx} = 1$, dont la somme vaut 0.

Pour $g = e^x$~: $u = e^{x + (T - t)/2}$. Vérification~: $u_t = -\frac12u$ et $\frac12u_{xx} = \frac12u$, dont la somme vaut également 0.

Les différences finies et l'échantillonnage concordent avec les deux solutions aux points testés. Un taux $r$ multiplie chacune par $e^{-r(T - t)}$, ce qu'on vérifie par substitution. La valeur terminale vaut $g(x)$, comme il se doit.

\iqwrong{Se tromper de signe sur $u_t$. C'est une équation rétrograde, résolue à partir de l'instant terminal. Le temps restant, $T - t$, apparaît donc, et non $t$.}

\iqfollow{Comment obtenez-vous l'équation de Black--Scholes à partir de ce résultat~?}{On utilise $dS = rS\,dt + \sigma S\,dW$ sous la probabilité risque-neutre. Alors la formule dit que le prix vérifie $u_t + rSu_S + \frac12\sigma^2S^2u_{SS} - ru = 0$, qui est cette équation.}

\iqwhy{C'est le lien entre espérances et équations différentielles, qui permet de choisir le plus simple des deux.}
\end{solution}

\begin{solution}{iq:iv1:stochastic-calculus:20}
$dX = X\,dW$. L'espérance vaut toujours 1, la médiane vaut $e^{-T/2}$, et $P(X_T < 1) = \Phi(\sqrt{T}/2)$, qui vaut $0.69$ en $T = 1$ et $0.9999997$ en $T = 100$. Presque toutes les trajectoires tendent vers 0.

$X$ est la martingale exponentielle, avec $dX = X\,dW$ d'après le lemme d'Itô. Donc l'espérance reste à 1.

$\ln X_T = W_T - T/2$ est gaussien d'espérance $-T/2$ et de variance $T$. Donc la médiane vaut $e^{-T/2}$~: $0.607$ en $T = 1$ et environ $2 \times 10^{-22}$ en $T = 100$. Par ailleurs, $X_T < 1$ exige $W_T < T/2$, dont la chance vaut $\Phi(\sqrt{T}/2)$\index{martingale exponentielle}.

Comme $W_t/t \to 0$, $\ln X_t \approx -t/2 \to -\infty$, donc chaque trajectoire tend vers 0 avec probabilité un.

Le rapprochement tient à ce que l'espérance est portée par des trajectoires de plus en plus rares, aux valeurs énormes. La variance vaut $e^T - 1$, ce qui est énorme. Pour $T = 100$, aucune simulation réalisable ne peut retrouver une espérance de 1. Un million de trajectoires donnent en général beaucoup moins, le reste se logeant dans une queue trop fine pour être échantillonnée. Notons que l'arrêt optionnel ne s'applique pas en $T = \infty$~: la limite vaut 0, alors que l'espérance vaut 1, si bien que la martingale n'est pas uniformément intégrable.

En termes d'investissement, il s'agit d'un pari équitable à pleine taille, avec une chance positive de gains énormes et une ruine presque sûre à long terme. C'est pourquoi on s'intéresse aux taux de croissance et pas seulement aux espérances.

\iqwrong{Affirmer que l'espérance est un bon guide de la richesse typique. C'est l'exemple extrême d'une espérance qui ne décrit aucune trajectoire qu'on ait des chances de voir.}

\iqfollow{Quels sont les moments d'ordre deux en $T = 1$ et $T = 100$~?}{$e^1 = 2.72$ et $e^{100} = 2.7 \times 10^{43}$. Donc la moyenne empirique est inutile bien avant $T = 100$.}

\iqwhy{Cela sépare les moyennes d'ensemble des moyennes temporelles, ce qui est au cœur de la discussion sur le dimensionnement et sur le critère de Kelly.}
\end{solution}

\begin{solution}{iq:iv1:stochastic-calculus:21}
L'espérance conditionnelle vaut $TW_T/2$, et la variance conditionnelle vaut $T^3/12$. Pour $T = 1$, elles valent $W_1/2$ et $\frac{1}{12}$, et la valeur finale explique $\frac34$ de la variance.

$I_T$ et $W_T$ sont conjointement gaussiens, centrés, avec $\operatorname{Var}I_T = T^3/3$, $\operatorname{Var}W_T = T$ et $\operatorname{Cov} = T^2/2$. Ces relations découlent de $I_T = \int_0^T(T - t)\,dW_t$.

La régression gaussienne\index{gaussienne conditionnelle} s'applique donc. La pente vaut $\operatorname{Cov}/\operatorname{Var}W_T = T/2$, ce qui donne l'espérance $TW_T/2$. Cela se comprend~: sachant la valeur finale, le pont brownien est en moyenne une droite allant de 0 à $W_T$, dont l'intégrale est l'aire d'un triangle.

La variance restante vaut
\[
\frac{T^3}{3} - \frac{(T^2/2)^2}{T} = \frac{T^3}{3} - \frac{T^3}{4} = \frac{T^3}{12}.
\]
Une seconde voie~: elle doit égaler la variance de l'intégrale d'un pont brownien sur $[0, T]$. Pour $T = 1$, $\int_0^1\int_0^1(\min(s, t) - st)\,ds\,dt = \frac13 - \frac14 = \frac{1}{12}$, soit la même valeur.

La part expliquée vaut $(\frac13 - \frac{1}{12})/\frac13 = \frac34$, soit $0.866^2$, le carré de la corrélation entre les deux. Une simulation à $40\,000$ trajectoires retrouve la pente, la variance résiduelle et la part.

\iqwrong{Croire que l'intégrale est indépendante de la valeur finale parce que les accroissements sont indépendants. Or elles ne sont pas indépendantes~: toutes deux sont des sommes des mêmes accroissements.}

\iqfollow{Cela fournit-il une variable de contrôle pour une option asiatique~?}{Oui. Comme la valeur finale explique les trois quarts de la variance, une moyenne simulée peut être régressée sur $W_T$, et cela seul réduit la variance de cette même part. Une moyenne géométrique fait encore mieux.}

\iqwhy{Conditionner sur une variable gaussienne est un outil d'une ligne, et ce même outil apparaît dans les constructions par pont brownien et dans la réduction de variance.}
\end{solution}

\begin{solution}{iq:iv1:stochastic-calculus:22}
$A$ suit la loi de l'arcsinus, de densité $1/(\pi\sqrt{x(1 - x)})$~: espérance $\frac12$, variance $\frac18$, et $P(A \le 0.1) = \frac{2}{\pi}\arcsin\sqrt{0.1} = 0.205$. Ainsi, un trader dépourvu de talent est en avance pendant 90~\% de l'année dans environ 20~\% des cas.

La densité est maximale près de 0 et de 1, et minimale en $\frac12$\index{loi de l'arcsinus}. L'espérance découle de la symétrie. La variance vient de $\mathbb{E}A^2 = \frac38$, une intégrale double de la probabilité d'orthant.

La fonction de répartition vaut $F(x) = \frac{2}{\pi}\arcsin\sqrt{x}$, donc $P(A \le 0.1) = 0.205$. Par symétrie, $P(A \ge 0.9)$ vaut autant. Ainsi, avec une chance de $41~\%$, la trajectoire passe plus de 90~\% du temps d'un seul côté. Il faut comparer une fraction comprise entre $0.45$ et $0.55$, qui n'a qu'une chance de $6.4~\%$. Le résultat le plus équilibré est le moins probable.

Une marche simulée à 2000 pas donne une espérance de $0.50$, une variance de $0.125$ et une masse en queue voisine de $0.205$, ce qui constitue une vérification indépendante de la formule. Retenons la conclusion. Une courbe de résultat positive pendant 90~\% de l'année en dit peu sur le talent, puisqu'environ un trader sur cinq, dénué de tout talent, en aura une.

En revanche, l'edge d'un trader fait tendre la fraction vers 1 de manière fiable, mais cette mesure est un test faible de celui-ci. Une statistique $t$ sur la moyenne exploite mieux les mêmes données.

\iqwrong{Supposer qu'une trajectoire d'espérance nulle doit passer environ la moitié de son temps au-dessus de 0. C'est l'inverse qui est vrai~: de longs séjours d'un seul côté sont normaux, puisque les excursions sont longues.}

\iqfollow{Quel est l'écart type de $A$~?}{$\sqrt{1/8} = 0.354$. C'est au-dessus de l'écart type de la loi uniforme, $0.289$, et déjà bien en chemin vers le $0.5$ d'un mélange à parts égales de 0 et de 1.}

\iqwhy{C'est le résultat sur les marches aléatoires qui heurte le plus l'intuition, et il se lit directement dans l'évaluation des performances.}
\end{solution}

\begin{solution}{iq:iv1:stochastic-calculus:23}
$P(M_T \ge a, W_T \le x) = P(W_T \ge 2a - x)$. L'écart $M_T - W_T$ suit la loi de $|W_T|$. L'espérance de l'écart vaut 16~\%, et la chance qu'il dépasse 20~\% vaut $0.317$.

Pour $x \le a$, on réfléchit chaque trajectoire après sa première atteinte de $a$\index{principe de réflexion}. Une trajectoire avec $M_T \ge a$ et une fin inférieure ou égale à $x$ correspond de façon biunivoque à une trajectoire finissant supérieure ou égale à $2a - x$, qui a nécessairement dû franchir $a$ en chemin. Donc les deux événements ont la même probabilité.

En dérivant par rapport à $a$ et $x$, on obtient la densité conjointe de $(M_T, W_T)$. À partir de celle-ci, la densité de $M_T - W_T$ se calcule directement, et c'est celle de $|W_T|$.

Un argument plus net utilise le renversement du temps\index{renversement du temps}. Posons $\hat{W}_s = W_{T - s} - W_T$, qui est aussi un mouvement brownien. Alors $M_T - W_T = \max_s(W_s - W_T) = \max_s\hat{W}_{T - s} = \max_s\hat{W}_s$, dont la loi est celle du maximum, donc celle de $|W_T|$ par réflexion. D'où
\[
\mathbb{E}(M_T - W_T) = \sqrt{2T/\pi}, \qquad P(M_T - W_T \ge c) = 2(1 - \Phi(c/\sqrt{T})).
\]

Avec une volatilité $\sigma$, il suffit de multiplier par $\sigma$. Donc $0.798 \times 0.2 = 16~\%$, et $P(\text{écart} > 20~\%) = 2(1 - \Phi(1)) = 0.317$. Une simulation à $40\,000$ trajectoires concorde après correction de l'observation discrète. Comme $\mathbb{E}W_T = 0$, ce même nombre est aussi égal à $\mathbb{E}M_T$.

\iqwrong{Voir dans un écart de 16~\% par rapport au plus haut le signe que quelque chose cloche. Pour un processus sans dérive et de volatilité 20~\%, c'est à quoi ressemble une année moyenne.}

\iqfollow{Quelle est l'espérance de l'écart après deux ans~?}{Il croît avec la racine du temps~: $16~\% \times \sqrt{2} = 22.6~\%$. Notons qu'il s'agit de l'écart courant en fin de période. Le drawdown maximal sur la période est plus grand.}

\iqwhy{Le renversement du temps est un tour élégant, et il fournit un indice de référence pour les drawdowns, que les investisseurs estiment souvent mal.}
\end{solution}

\begin{solution}{iq:iv1:stochastic-calculus:24}
$\mathbb{E}\tau = a/\mu = 10$, et $\operatorname{Var}\tau = a\sigma^2/\mu^3 = 40$. $\tau$ suit la loi gaussienne inverse, avec $P(\tau > 20) = 0.073$.

$\tau$ a une espérance finie sous $\mu > 0$. Pour l'espérance, $X_t - \mu t$ est une martingale. Le théorème d'arrêt optionnel\index{théorème d'arrêt optionnel} donne $a - \mu\mathbb{E}\tau = 0$.

Pour la variance, $(X_t - \mu t)^2 - \sigma^2t$ est une martingale. En $\tau$, $X_\tau = a$, donc
\[
\mathbb{E}(a - \mu\tau)^2 = \sigma^2\mathbb{E}\tau.
\]
Le membre de gauche vaut $\mu^2\operatorname{Var}\tau$, puisque $\mathbb{E}\mu\tau = a$. Donc $\operatorname{Var}\tau = \sigma^2(a/\mu)/\mu^2 = a\sigma^2/\mu^3$.

Pour les valeurs données~: $\mathbb{E}\tau = 1/0.1 = 10$, et $\operatorname{Var}\tau = 0.04/0.001 = 40$, avec un écart type de $6.3$.

Il s'agit de la loi gaussienne inverse\index{loi gaussienne inverse}, de moyenne $a/\mu$ et de paramètre de forme $a^2/\sigma^2 = 25$. Sa queue admet la forme fermée
\[
P(\tau \le t) = \Phi\left(\frac{\mu t - a}{\sigma\sqrt{t}}\right) + e^{2\mu a/\sigma^2}\Phi\left(\frac{-\mu t - a}{\sigma\sqrt{t}}\right),
\]
qui donne $0.073$ pour $P(\tau > 20)$. Un tirage de la loi et une simulation d'Euler directe du processus concordent avec la moyenne et la variance.

Limites~: quand $\mu \to 0$, les deux explosent, et la variance plus vite, comme il se doit, puisque le cas sans dérive a une espérance infinie. Pour $\mu = 0.01$, la variance atteint déjà $40\,000$.

\iqwrong{Donner le temps déterministe $a/\mu$ sans dispersion, ou une variance de $\sigma^2/\mu^2$, qui a la mauvaise échelle en $a$ et en $\mu$. L'analyse dimensionnelle permet de le repérer.}

\iqfollow{Si la cible est 10 fois plus éloignée, que devient la dispersion relative~?}{L'espérance et la variance sont multipliées chacune par 10, donc le coefficient de variation diminue de $\sqrt{10}$. Les cibles à long horizon sont plus prévisibles, en proportion.}

\iqwhy{Le temps d'atteinte d'une cible est courant en exécution et en évaluation de stratégies. La démonstration par martingale va plus vite que la densité.}
\end{solution}

\begin{solution}{iq:iv1:stochastic-calculus:25}
Avec $\lambda = \theta^2/2$, la martingale $e^{\theta W_t - \theta^2t/2}$ donne le résultat. La limite $\lambda \to 0$ montre que $T_a$ est fini, et la dérivée infinie montre que son espérance est infinie. Pour $a = 1$, $\lambda = \frac12$, la valeur vaut $e^{-1} = 0.368$.

Pour $\theta > 0$, $M_t = e^{\theta W_t - \theta^2t/2}$ est une martingale\index{martingale exponentielle}. Arrêtée en $T_a$, elle est majorée par $e^{\theta a}$, puisque $W \le a$ avant $T_a$ et que le terme en temps ne fait que la réduire. Sur l'événement $T_a = \infty$, $M_t \to 0$, car $W_t$ reste sous $a$ pendant que $t$ croît. Le théorème d'arrêt optionnel borné donne alors
\[
1 = \mathbb{E}\left[e^{\theta a}e^{-\theta^2T_a/2}\mathbf{1}(T_a < \infty)\right].
\]
En posant $\lambda = \theta^2/2$, on a $\mathbb{E}[e^{-\lambda T_a}\mathbf{1}(T_a < \infty)] = e^{-a\sqrt{2\lambda}}$.

La racine positive compte. Avec $\theta < 0$, $M$ n'est pas bornée lorsque $W \to -\infty$ avant $T_a$, donc le théorème d'arrêt échoue. L'autre racine donnerait une transformée supérieure à 1, ce qui est impossible.

Lorsque $\lambda \to 0$~: $P(T_a < \infty) = 1$. La dérivée de $e^{-a\sqrt{2\lambda}}$ en 0 vaut $-\infty$, à cause de la racine carrée, et donc $\mathbb{E}T_a = \infty$.

Pour $a = 1$, $\lambda = \frac12$~: $e^{-1}$. Un tirage de $T_1 = 1/Z^2$ donne $0.368$.

\iqwrong{Appliquer le théorème d'arrêt optionnel sans vérifier la bornitude. La même étape montrerait que $\mathbb{E}W_{T_a} = 0$, ce qui est faux, puisque $W_{T_a} = a$.}

\iqfollow{Que change une dérive $\mu > 0$~?}{Le même argument avec un exposant décalé donne $\mathbb{E}e^{-\lambda\tau} = e^{a(\mu - \sqrt{\mu^2 + 2\lambda})}$, avec $\sigma = 1$. La dérivée en 0 est maintenant finie, égale à $-a/\mu$. Pour $\mu = \frac12$, $a = 1$, $\lambda = \frac12$, cela donne $e^{0.5 - 1.118}$, comme le confirme un tirage.}

\iqwhy{C'est un usage compact du théorème d'arrêt, et le choix de la racine montre si ses hypothèses sont comprises.}
\end{solution}

\begin{solution}{iq:iv1:stochastic-calculus:26}
À $\mu = 10~\%$, le temps moyen vaut $\ln 2/0.08 = 8.66$ ans. À $\mu = 1~\%$, la probabilité vaut $2^{-1/2} = 0.707$. À $\mu = 2~\%$, le processus atteint 2 avec certitude, mais le temps moyen est infini.

$X = \ln S$ a pour dérive $\nu = \mu - \frac12\sigma^2$ et pour volatilité $\sigma$, et la cible est $a = \ln 2$\index{temps de premier passage}.

$\mu = 10~\%$~: $\nu = 8~\%$. La moyenne d'un temps d'atteinte avec dérive positive vaut $a/\nu = 0.693/0.08 = 8.66$.

$\mu = 1~\%$~: $\nu = -1~\%$, donc le logarithme s'éloigne. La probabilité d'atteindre un jour le niveau $a$ situé au-dessus vaut $e^{2\nu a/\sigma^2} = e^{-0.5\ln 2} = 2^{-1/2}$. Notons que l'espérance de $S$ continue de croître.

Vérification par martingale~: $p = 1 - 2\mu/\sigma^2 = \frac12$, donc $\sqrt{S}$ est une martingale\index{fonction d'échelle}. Elle est bornée par $\sqrt{2}$ jusqu'au temps d'atteinte, et si 2 n'est jamais atteint, $S \to 0$ puisque $\nu < 0$, et la martingale tend vers 0. Donc $1 = q\sqrt{2}$, ce qui donne $q = 1/\sqrt{2}$. Les réponses concordent, sans recours à la première formule.

$\mu = 2~\%$~: $\nu = 0$. Le logarithme est sans dérive, il atteint donc la cible avec certitude et a une espérance infinie, comme c'est le cas pour l'atteinte d'un seul côté par un mouvement brownien sans dérive.

Une simulation sur 400 ans à $\mu = 1~\%$ atteint 2 dans environ $68~\%$ des trajectoires, après correction des observations discrètes. Cela concorde avec la formule du premier passage à cet horizon, $0.683$, et la limite vaut $0.707$.

\iqwrong{Affirmer qu'une $\mu$ positive signifie que l'actif double à coup sûr. Le signe qui compte est celui de la dérive du logarithme. Avec $0 < \mu < \sigma^2/2$, l'espérance monte alors que la trajectoire typique baisse.}

\iqfollow{À partir de quelle valeur de $\mu$ le temps moyen devient-il fini~?}{Pour toute $\mu > \sigma^2/2 = 2~\%$, et il explose à cette valeur. Le temps moyen pour $\mu = 2.1~\%$ serait d'environ 693 ans.}

\iqwhy{Il teste la différence entre dérive arithmétique et dérive géométrique, et la vérification par une martingale puissance est une bonne habitude.}
\end{solution}

\begin{solution}{iq:iv1:stochastic-calculus:27}
$|W_t| = \int_0^t\operatorname{sign}(W_s)\,dW_s + L_t$, avec $\mathbb{E}L_1 = \sqrt{2/\pi} = 0.798$. $L$ mesure le temps passé en 0, par unité d'espace, et il a la loi du maximum courant.

Une approximation lisse de $|x|$ a une dérivée seconde concentrée autour de 0, de masse totale 2. Donc le terme d'Itô $\frac12f^{\prime\prime}(W)\,dt$ ne s'annule pas, et il ne s'accumule que lorsque $W$ est en 0\index{formule de Tanaka}. Sa limite est $L_t$, définie par
\[
L_t = \lim_{\varepsilon \to 0}\frac{1}{2\varepsilon}\int_0^t\mathbf{1}(|W_s| < \varepsilon)\,ds.
\]
C'est le temps local\index{temps local} en 0~: un taux de temps passé près de 0 par unité de largeur.

L'intégrale stochastique est d'espérance nulle, donc $\mathbb{E}L_1 = \mathbb{E}|W_1| = \sqrt{2/\pi}$.

Il existe deux estimateurs à partir d'une trajectoire. Le premier est la formule d'occupation, avec un $\varepsilon$ petit. L'autre est $|W_1| - \sum\operatorname{sign}(W_{i})\Delta W_i$, qui réarrange la formule de Tanaka. Dans une simulation, chacun a une moyenne voisine de $0.798$, et les deux estimateurs sont fortement corrélés d'une trajectoire à l'autre, au-delà de $0.9$. Ils mesurent donc la même chose.

Par Tanaka, $\int\operatorname{sign}(W)\,dW$ est elle-même un mouvement brownien $B$, par la caractérisation de Lévy, puisque sa variation quadratique vaut $t$. En écrivant $|W_t| = B_t + L_t$ et en remarquant que $L$ ne croît que lorsque $|W| = 0$, le lemme de Skorokhod donne $L_t = \max_{s \le t}(-B_s)$. Donc $L_t$ a la loi du maximum d'un mouvement brownien, et donc celle de $|W_t|$. Les quantiles de la simulation concordent.

\iqwrong{Appliquer le lemme d'Itô à $|x|$ avec $f^{\prime\prime} = 0$, ce qui donne $|W_t| = \int\operatorname{sign}(W)\,dW$. Son espérance serait nulle, ce qui est impossible pour une variable positive.}

\iqfollow{Où cela apparaît-il en couverture~?}{Une stratégie de stop-loss ne détient un actif que lorsque le cours est au-dessus du strike. Son échec à répliquer un call correspond exactement au temps local au strike~: un coût qui ne disparaît pas avec des transactions plus fréquentes.}

\iqwhy{C'est l'un des rares endroits où le terme de dérivée seconde du lemme d'Itô apparaît seul. Il met à l'épreuve les candidats confirmés, et le lien avec la couverture est réel.}
\end{solution}

\begin{solution}{iq:iv1:stochastic-calculus:28}
$\mathbb{E}Q = \frac12$, et $\operatorname{Var}Q = \frac13$. La série donne les mêmes résultats, en utilisant aussi $\sum_k((k - \frac12)\pi)^{-4} = \frac16$.

L'espérance vaut $\int_0^1\mathbb{E}W_t^2\,dt = \int_0^1t\,dt = \frac12$.

Pour la variance, $\operatorname{Cov}(W_s^2, W_t^2) = 2\min(s, t)^2$, comme pour deux gaussiennes quelconques de moyenne nulle. Donc
\[
\operatorname{Var}Q = 2\int_0^1\int_0^1\min(s, t)^2\,ds\,dt = 4\int_0^1\int_0^ts^2\,ds\,dt = 4\int_0^1\frac{t^3}{3}\,dt = \frac13.
\]

Pour la série, les fonctions sinus $\sqrt{2}\sin((k - \frac12)\pi t)$ sont orthonormées sur $[0, 1]$\index{Karhunen-Loève, développement de}. En posant $\lambda_k = ((k - \frac12)\pi)^{-2}$ et en intégrant le carré du développement,
\[
Q = \sum_k\lambda_kZ_k^2,
\]
les termes croisés s'annulant. Donc $\mathbb{E}Q = \sum_k\lambda_k = \frac12$, comme annoncé, et $\operatorname{Var}Q = 2\sum_k\lambda_k^2$, puisque $\operatorname{Var}Z^2 = 2$. La somme demandée est donc $\sum_k\lambda_k^2 = \frac16$. C'est la somme classique $\sum 1/(2k - 1)^4 = \pi^4/96$ déguisée, et une somme numérique donne $\frac16$ à douze décimales. Les deux voies concordent donc.

Le premier terme a $\lambda_1 = 4/\pi^2 = 0.405$ et porte $98.6~\%$ de la variance. Donc $Q$ est proche d'un chi-deux à un degré de liberté à l'échelle près, et sa loi est fortement dissymétrique. Une simulation donne $0.50$ et $0.33$.

\iqwrong{En utilisant $\mathbb{E}Q^2 = (\mathbb{E}Q)^2 + \int\operatorname{Var}W_t^2\,dt$, ce qui ignore la corrélation dans le temps. On obtient $\frac14 + \frac23$ au lieu de $\frac14 + \frac13$, ce qui est trop grand.}

\iqfollow{Pourquoi un seul terme porte-t-il à lui seul autant de la variance~?}{Les valeurs propres décroissent en $1/k^2$, et leurs carrés en $1/k^4$. Le niveau d'une marche aléatoire est dominé par sa composante la plus lente.}

\iqwhy{De telles quantités apparaissent en statistique des racines unitaires, dont les limites sont des intégrales comme $Q$. Le développement en dit bien plus que les deux premiers moments.}
\end{solution}

\begin{solution}{iq:iv1:stochastic-calculus:29}
$u(0, 0) = 1/\sqrt{\cosh\theta}$, avec $A = \theta\tanh(\theta(1 - t))$, et $B$ son intégrale. Pour $\theta = 1$, la valeur vaut $0.805$.

On pose $u(t, x) = \mathbb{E}[e^{-\frac12\theta^2\int_t^1W_s^2\,ds} \mid W_t = x]$. La formule de Feynman--Kac\index{formule de Feynman-Kac} avec un facteur d'actualisation dépendant de l'état donne
\[
u_t + \frac12u_{xx} - \frac12\theta^2x^2u = 0, \qquad u(1, x) = 1.
\]
On substitue $u = \exp(-\frac12A(t)x^2 + B(t))$. Alors $u_t = (-\frac12A'x^2 + B')u$ et $u_{xx} = (A^2x^2 - A)u$. En identifiant les puissances de $x$~:
\[
A' = A^2 - \theta^2, \qquad B' = \tfrac12A, \qquad A(1) = B(1) = 0.
\]
L'équation de Riccati\index{equation de Riccati@équation de Riccati} admet pour solution $A(t) = \theta\tanh(\theta(1 - t))$. Vérification~: $A' = -\theta^2\operatorname{sech}^2 = \theta^2(\tanh^2 - 1) = A^2 - \theta^2$.

Alors
\[
B(0) = -\int_0^1\frac12\theta\tanh(\theta(1 - t))\,dt = -\frac12\ln\cosh\theta.
\]
Donc $u(0, 0) = e^{B(0)} = (\cosh\theta)^{-1/2}$\index{formule de Cameron-Martin}.

Pour $\theta$ petit, cela vaut $1 - \theta^2/4$, tandis que le membre de gauche se développe en $1 - \frac12\theta^2\mathbb{E}Q = 1 - \theta^2/4$. Les deux concordent. Quand $\theta \to \infty$, cela décroît comme $\sqrt{2}e^{-\theta/2}$, si bien que la probabilité d'un $Q$ très petit chute rapidement, comme il se doit.

Une seconde voie passe par la série~: $Q = \sum\lambda_kZ_k^2$ donne $\prod_k(1 + \theta^2\lambda_k)^{-1/2}$, et la forme en produit du cosinus hyperbolique ramène au même résultat. Un produit numérique et une simulation donnent tous deux $0.805$.

\iqwrong{En essayant une forme sans terme en $x^2$, ou en se trompant de signe dans le $\tanh$ de sorte que $u$ croisse avec $x$. L'actualisation pénalise les grands $x$, donc $u$ doit décroître avec lui.}

\iqfollow{Pourquoi un exposant quadratique est-il la bonne hypothèse~?}{Le potentiel est quadratique et la diffusion est gaussienne. Ces problèmes restent gaussiens, et un exposant quadratique est la forme gaussienne la plus générale.}

\iqwhy{Cela demande plusieurs techniques à la fois. La structure de Riccati est la même que celle qui sous-tend les modèles quadratiques-gaussiens et affines.}
\end{solution}

\begin{solution}{iq:iv1:stochastic-calculus:30}
$\mathbb{E}e^{-\lambda\tau} = 1/\cosh(a\sqrt{2\lambda})$, $\mathbb{E}\tau = a^2$, $\operatorname{Var}\tau = \frac23a^4$, et $P(\tau > t)$ décroît au taux $\pi^2/(8a^2)$.

$\cosh(\theta W_t)e^{-\theta^2t/2}$ est une martingale, comme moyenne de $e^{\pm\theta W_t - \theta^2t/2}$\index{temps de sortie}. Ce processus est borné jusqu'à $\tau$, et en $\tau$ on a $|W_\tau| = a$, donc $\cosh(\theta W_\tau) = \cosh\theta a$. Le théorème d'arrêt optionnel donne $\cosh(\theta a)\mathbb{E}e^{-\theta^2\tau/2} = 1$, et avec $\lambda = \theta^2/2$ on obtient le résultat.

En développant avec $x = a\sqrt{2\lambda}$~: $1/\cosh x = 1 - \frac12x^2 + \frac{5}{24}x^4 - \dots$, et $x^2 = 2\lambda a^2$. Donc
\[
\mathbb{E}e^{-\lambda\tau} = 1 - \lambda a^2 + \frac56\lambda^2a^4 - \dots
\]
En identifiant avec $1 - \lambda\mathbb{E}\tau + \frac12\lambda^2\mathbb{E}\tau^2$~: $\mathbb{E}\tau = a^2$, et $\mathbb{E}\tau^2 = \frac53a^4$. Donc $\operatorname{Var}\tau = \frac23a^4$.

Vérification~: le polynôme est une martingale\index{martingales de Hermite}. À $\tau$, $a^4 - 6a^2\mathbb{E}\tau + 3\mathbb{E}\tau^2 = 0$. Avec $\mathbb{E}\tau = a^2$, cela donne $\mathbb{E}\tau^2 = \frac53a^4$, ce qui redonne le même résultat.

Pour la queue, $\cosh(a\sqrt{2\lambda})$ s'annule lorsque $a\sqrt{2\lambda} = i\pi/2$, donc en $\lambda = -\pi^2/(8a^2)$. La singularité la plus proche fixe la décroissance exponentielle de la queue\index{décroissance de la queue}~: $P(\tau > t) \approx \frac{4}{\pi}e^{-\pi^2t/(8a^2)}$.

Pour $a = 1$, $\lambda = \frac12$~: $1/\cosh 1 = 0.648$, et $P(\tau > 4) \approx 0.0092$. Une marche symétrique sur réseau à $100\,000$ trajectoires s'accorde sur la moyenne, la variance, la transformée et la queue. Par comparaison avec le cas à une seule barrière, une seconde barrière transforme la queue lourde, d'ordre $1/\sqrt{t}$, en une queue exponentielle, et tous les moments deviennent finis.

\iqwrong{En utilisant $\mathbb{E}\tau^2 = (\mathbb{E}\tau)^2 + \mathbb{E}\tau$, ou en omettant le facteur $\frac12$ dans le développement de l'exponentielle. La martingale polynomiale permet de détecter rapidement ce type d'erreur.}

\begin{sloppypar}\iqfollow{Pourquoi $\frac{4}{\pi}$ est-il la constante de la queue~?}{La queue complète est la série $\frac{4}{\pi}\sum_{k \ge 0}\frac{(-1)^k}{2k + 1}e^{-(2k + 1)^2\pi^2t/(8a^2)}$, et le premier terme domine rapidement. En $t = 0$, la somme alternée vaut $\pi/4$, ce qui donne 1, comme il se doit.}\end{sloppypar}

\iqwhy{Il relie une martingale, une transformée, des moments et une queue dans un même problème. La vérification par une seconde martingale polynomiale illustre le réflexe attendu au plus haut niveau.}
\end{solution}
